\section{The approximation threshold}
\label{sec:conclusion}

The construction and soundness analysis were stated with explicit
requirements on their constants. We now choose those constants jointly,
before requesting the source alphabets.

For the given integer \(m\geq4\), the definitions in
\Cref{src:construction-parameters,src:variance-parameters,src:list-parameters,src:decoding-soundness}
are completed by choosing \(d=2^{38}m^{10}\). Thus
\begin{equation}
\begin{aligned}
 a&=\frac1{4m},& b_0&=\frac1{2m},
 &\nu&=\frac1{2^7m^3},\\
 \beta&=\frac1{2^{10}m^3},&\ell&=2^{11}m^3,
 &h&=2^{12}m^3,\\
 d&=2^{38}m^{10},&\sigma&=\frac1{2^{46}m^{12}},
 &p&=2d+1,\\
 t&=2^{17}m^4.
\end{aligned}
\label{src:parameters}
\end{equation}
The construction conditions \eqref{src:dimension-bounds} hold because
\(d\) is even and \(\sqrt d=2^{19}m^5>4m\). For the restriction
conditions \eqref{src:restriction-budget}, we have \(h<d\) and
\[
 16h+32h^2=2^{16}m^3+2^{29}m^6
 \leq2^{30}m^7=\frac{\nu d}{2}.
\]
The displayed inequality holds already for \(m\geq1\).
The remaining list and decoding identities are
\[
 \frac{\beta h}{2}=2,\qquad
 \binom h2\ell^2\sigma=\frac{h-1}{2h}<\frac12.
\]
All choices depend only on \(m\); no label alphabet has entered them.
Their magnitude affects the constants and degree of the polynomial
running time, but every fixed \(m\) gives a fixed reduction.

\begin{proof}[Proof of \Cref{thm:gap,thm:hardness}]
Fix \(m\geq4\) and the constants in \Cref{src:parameters}.
Apply \Cref{src:hardness} at this \(\sigma\) to obtain fixed
source alphabets and a polynomial-time reduction from \(\mathrm{3SAT}\).
Now define \(q,r\) from those alphabets and construct the graph.
\Cref{graph:polynomial} makes its vertex and edge lists explicit in
deterministic polynomial time. All the hypotheses of
\Cref{graph:completeness,dec:soundness} have been verified above, so
these results give the two cover bounds in \Cref{thm:gap}.

For a fixed real \(\alpha\in[1,2)\), choose a fixed integer
\(m\geq4\) such that
\begin{equation}
 \alpha<R_m:=\frac{1-1/m}{1/2+1/m}
       =2-\frac6{m+2}.
\label{eq:gap-ratio}
\end{equation}
Suppose an algorithm always returns a vertex cover of size at most
\(\alpha\tau(G)\). On a completeness instance its output has
size strictly less than
\[
 \alpha\left(\frac12+\frac1m\right)n
 <\left(1-\frac1m\right)n,
\]
whereas on a soundness instance every cover has size strictly greater
than \((1-1/m)n\). Comparing the output cardinality with this
rational threshold distinguishes the two source cases. The comparison
uses only integer arithmetic: multiply the cardinality by \(m\)
and compare with \((m-1)n\). The real number \(\alpha\) is used
only to choose the fixed integer \(m\); it is not an input to the
reduction. This proves \Cref{thm:hardness}.
\end{proof}

For completeness, the matching upper bound on the approximation
factor is elementary. Find a maximal matching and return the set of
all its endpoints. Maximality makes this set a vertex cover: an edge
with neither endpoint selected could be added to the matching.
Every vertex cover contains an endpoint of each matching edge, and
these edges are disjoint. The returned cover therefore has size at
most twice optimum. Together with \Cref{thm:hardness}, this gives the
constant approximation threshold \(2\), unless
\(\mathrm P=\mathrm{NP}\).
